\documentclass[11pt]{article}

\usepackage[T1]{fontenc}
\usepackage[utf8]{inputenc}
\usepackage[spanish,es-nodecimaldot]{babel}
\usepackage{amsmath}
\usepackage{amssymb}
\usepackage{booktabs}
\usepackage{geometry}
\usepackage{hyperref}
\usepackage{xcolor}

\geometry{margin=2.5cm}
\hypersetup{
  colorlinks=true,
  linkcolor=blue!55!black,
  urlcolor=blue!55!black
}

\title{Función de recompensa de VSSS Lab\\
  \large MAPPO M24.3 con formación por roles dinámicos}
\author{VSSS Lab}
\date{3 de agosto de 2026}

\newcommand{\indicator}[1]{\mathbb{I}\!\left[#1\right]}
\newcommand{\Rteam}{R_t^{\mathrm{team}}}

\begin{document}

\maketitle

\begin{abstract}
Este documento describe la función de recompensa utilizada por la configuración
\texttt{m24-3-mappo-role-formation.toml}. La recompensa terminal continúa siendo global para
los tres robots. El cambio respecto al run 0013 añade crédito denso, acotado e independiente de
la identidad física para mejorar la formación de los roles dinámicos \emph{support} y
\emph{coverage}. El objetivo es reducir la persecución colectiva del balón sin crear un portero
fijo ni premiar al robot que da el último toque.
\end{abstract}

\section{Principio de distribución}

El entorno produce un escalar de equipo por mundo y por decisión. El colector lo replica para
los tres agentes activos:

\begin{equation}
  R_{t,0}=R_{t,1}=R_{t,2}=\Rteam.
\end{equation}

Por tanto, un gol propio paga $+10$ a cada agente y un gol recibido paga $-10$ a cada agente.
La recompensa no identifica al anotador. Premiar únicamente el último toque introduciría una
competencia directa por rematar y reforzaría precisamente la conducta de todos contra el balón.

Los roles \emph{attacker}, \emph{support} y \emph{coverage} se reasignan en cada estado mediante
una permutación de costo mínimo con histéresis. Ningún ID, marcador ni posición permanente del
arreglo es dueño de un rol.

\section{Ecuación completa}

La recompensa de equipo es la suma de términos nombrados:

\begin{align}
\Rteam ={}& R_{\mathrm{goal}} + R_{\mathrm{semantic}}
 + R_{\mathrm{mouth}} + R_{\mathrm{geometry}} + R_{\mathrm{formation}}
 + R_{\mathrm{touch}} + R_{\mathrm{defense}} \notag\\
&+ R_{\mathrm{time}} + R_{\mathrm{action}}
 + R_{\mathrm{congestion}} + R_{\mathrm{deadlock}}
 + R_{\mathrm{spin}} + R_{\mathrm{draw}} + R_{\mathrm{stagnation}}.
\end{align}

La configuración activa conserva $\gamma=0.99$ y los coeficientes de la
Tabla~\ref{tab:coefficients}.

\begin{table}[ht]
  \centering
  \caption{Coeficientes activos de recompensa.}
  \label{tab:coefficients}
  \begin{tabular}{@{}lll@{}}
    \toprule
    Término & Coeficiente & Interpretación \\
    \midrule
    Gol propio / recibido & $\pm 10$ & Resultado terminal dominante \\
    Resultado semántico & $\pm 1$ & Éxito, fracaso o timeout del ejercicio \\
    Transporte a boca de gol & $5.0$ & Cambio de potencial angular del balón \\
    Geometría atacante & $0.05$ & Shaping descontado attacker--balón--gol \\
    Formación dinámica & $0.10$ & Shaping de support y coverage activos \\
    Contacto útil & $0.15$ & Impulso firmado al entrar en contacto \\
    Cobertura defensiva & $0.10$ & Mejora al cubrir una amenaza \\
    Cambio de acción & $-0.002$ & Regularización cuadrática \\
    Congestión & $-0.001$ & Agrupamiento persistente \\
    Deadlock aliado & $-0.002$ & Contacto aliado sostenido \\
    Deadlock rival & $-0.0005$ & Empuje rival sin progreso \\
    Giro inútil & $-0.002$ & Giro lento, lejos del balón, tras gracia \\
    Empate & $-0.50$ & Resultado no resuelto \\
    Estancamiento & $-1.0$ & Terminal heredado; free-ball lo evita en juego \\
    Tiempo & $-0.05/1500$ & Costo constante por decisión \\
    \bottomrule
  \end{tabular}
\end{table}

Los términos históricos de progreso ingenuo, dirección horizontal del balón, alineación aislada
y esfuerzo de ruedas permanecen configurados en cero.

\section{Términos terminales y semánticos}

El resultado futbolístico es

\begin{equation}
R_{\mathrm{goal}} =
10\indicator{\text{gol propio}} - 10\indicator{\text{gol rival}}.
\end{equation}

En un ejercicio semántico se añade

\begin{equation}
R_{\mathrm{semantic}} =
\begin{cases}
 +1, & \text{éxito},\\
 -1, & \text{fracaso},\\
 -1, & \text{timeout no resuelto},\\
 0,  & \text{de otro modo}.
\end{cases}
\end{equation}

El timeout se penaliza para que abstenerse de intentar el ejercicio no sea mejor que fallarlo.
Esta señal es independiente del gol base: un ejercicio de disparo puede producir tanto el evento
físico de gol como su resultado semántico.

\section{Transporte del balón hacia una oportunidad convertible}

Sea $\Phi_{\mathrm{mouth}}(s)$ la anchura angular de la boca de gol observada desde el balón,
normalizada a $[0,1]$. Es máxima frente a la portería y colapsa cerca de una incidencia rasante,
por lo que llevar el balón a una esquina no cobra como si fuera una ocasión:

\begin{equation}
R_{\mathrm{mouth}} = 5
\left(\Phi_{\mathrm{mouth}}(s_{t+1})-\Phi_{\mathrm{mouth}}(s_t)\right).
\end{equation}

Este término es deliberadamente no descontado en su propia diferencia. Mantener el balón quieto
no genera recompensa y moverlo de una posición buena a una peor devuelve lo cobrado. Su
contribución máxima por episodio está acotada por $5$, frente a un gol que vale $10$.

\section{Geometría atacante}

El potencial atacante combina cuatro cantidades normalizadas:

\begin{equation}
\Phi_{\mathrm{geometry}} =
0.45 A + 0.25 O + 0.15 C + 0.15 P,
\end{equation}

donde $A$ es la alineación attacker--balón--gol, $O$ es el margen dentro de la apertura
utilizable, $C$ es la proximidad controlable y $P$ es el progreso longitudinal. Su shaping es

\begin{equation}
R_{\mathrm{geometry}} = 0.05
\left(0.99\Phi_{\mathrm{geometry}}(s_{t+1})-
\Phi_{\mathrm{geometry}}(s_t)\right).
\end{equation}

El potencial siguiente se fija en cero al terminar el episodio. Así no se paga simplemente por
finalizar en una pose favorable.

\section{Formación por roles dinámicos}

El nuevo término considera únicamente los robots activos que actualmente poseen las
responsabilidades \emph{support} y \emph{coverage}. Para un rol activo $r$, sea $d_r$ la distancia
entre el robot asignado y el objetivo geométrico ya empleado por el asignador de roles:

\begin{align}
q_r &= \exp\!\left(-\frac{d_r}{0.25\ \mathrm{m}}\right),\\
\Phi_{\mathrm{formation}} &=
\left(\prod_{r\in\mathcal{R}_{\mathrm{active}}}q_r\right)^{
1/|\mathcal{R}_{\mathrm{active}}|},
\quad \mathcal{R}_{\mathrm{active}}\subseteq
\{\mathrm{support},\mathrm{coverage}\}.
\end{align}

Si ninguno de esos roles está activo, como en un ejercicio 1v0, el potencial es cero. La
recompensa es

\begin{equation}
R_{\mathrm{formation}} = 0.10
\left(0.99\Phi_{\mathrm{formation}}(s_{t+1})-
\Phi_{\mathrm{formation}}(s_t)\right).
\end{equation}

El potencial siguiente también es cero en terminales. Quedarse estacionado no produce una renta
positiva. Cuando dos robots intercambian responsabilidades, el shaping sigue al rol y no a la
identidad física. No existe un action mask: coverage todavía puede desafiar el balón cuando la
asignación dinámica considera que debe rotar.

La media geométrica hace que support y coverage sean un cuello de botella conjunto: una
responsabilidad bien cubierta ya no puede ocultar el abandono de la otra. El coeficiente $0.10$
acota su contribución al retorno descontado a uno por ciento del valor de un gol. Es la mitad del
primer experimento de ADR 0023, que produjo separación parcial acompañada de exceso de
\texttt{STOP}.

\section{Contacto y defensa}

El contacto útil es un impulso firmado y sólo ocurre en el flanco de entrada del contacto:

\begin{equation}
R_{\mathrm{touch}} = 0.15\tanh\!\left(
\operatorname{sign}_{\mathrm{attack}}\Delta v_{x,\mathrm{ball}}
\right).
\end{equation}

El signo hace que impulsar el balón en sentido incorrecto cueste lo mismo que el avance
equivalente. El flanco evita cobrar repetidamente durante contacto sostenido.

La cobertura defensiva paga reducir la distancia $d_{\mathrm{def}}$ al punto de cobertura sólo
cuando existe amenaza $T\in[0,1]$:

\begin{equation}
R_{\mathrm{defense}} = 0.10T
\left(d_{\mathrm{def},t}-d_{\mathrm{def},t+1}\right).
\end{equation}

Los deadlocks se penalizan después de una gracia: contacto aliado sostenido cuando no se disputa
el balón y contacto rival sostenido cuando el balón permanece estancado.

\section{Regularización}

El cambio brusco de comandos se penaliza mediante

\begin{equation}
R_{\mathrm{action}}=-0.002\operatorname{mean}\!\left((a_t-a_{t-1})^2\right).
\end{equation}

La congestión y el giro inútil tienen coeficientes pequeños. El giro sólo se activa después de
$0.5$ segundos cuando el robot está lento, gira por encima de $1\ \mathrm{rad/s}$ y está fuera
de una envolvente de $0.12$ m alrededor del balón. Estas condiciones evitan castigar una
orientación legítima antes de un contacto.

\section{Evidencia y criterios para el siguiente run}

En el replay de la iteración 2000 del run 0013 se midió:

\begin{table}[ht]
  \centering
  \caption{Conducta por rol antes del nuevo shaping.}
  \begin{tabular}{@{}lrr@{}}
    \toprule
    Rol & Fracción \texttt{STRIKE} & Distancia media al balón \\
    \midrule
    attacker & $0.875$ & $0.194$ m \\
    support  & $0.551$ & $0.434$ m \\
    coverage & $0.482$ & $0.455$ m \\
    \bottomrule
  \end{tabular}
\end{table}

La diferenciación espacial existía, pero support y coverage seguían solicitando contacto con el
balón aproximadamente la mitad del tiempo. El siguiente run debe evaluarse por conducta y
resultado, no por recompensa total. Los criterios principales son:

\begin{itemize}
  \item reducir \texttt{coverage.strike\_fraction} por debajo de $0.482$ sin colapsar a
    \texttt{STOP};
  \item reducir \texttt{support.strike\_fraction} por debajo de $0.551$ sin degradar pase o
    rotación;
  \item conservar attacker como el rol con mayor fracción de \texttt{STRIKE};
  \item mantener al menos $0.2$ goles por minuto y una tasa de empate no mayor que $0.70$;
  \item conservar la ocupación de posiciones convertibles que justificó el coeficiente de
    transporte $5.0$.
\end{itemize}

\section{Trazabilidad y rollback}

La implementación vive en \texttt{python/vsss\_train/marl\_env.py}; los coeficientes del nuevo
experimento están en \texttt{experiments/configs/m24-3-mappo-role-formation.toml}. La decisión
arquitectónica original es ADR 0023 y su corrección posterior al run role-0001 es ADR 0024. La
evidencia se registra en
\texttt{docs/evidence/m24-3-role-formation-credit.md}.

El rollback no requiere cambiar código: fijar
\texttt{role\_formation\_coefficient = 0.0} desactiva el término. La configuración histórica
\texttt{m24-3-mappo-circular.toml} permanece sin modificar y conserva el fingerprint del run
0013.

\end{document}
